\documentclass[../main.tex]{subfiles}
\graphicspath{{\subfix{../images/}}}
\begin{document}

\usetikzlibrary{arrows.meta,calc,positioning}

% Gleiche Farben wie in der Mindmap
\definecolor{bshaupt}{RGB}{40,90,188}
\definecolor{bshell}{RGB}{232,241,250}

\begin{figure}[htbp]
    \centering
    \resizebox{\textwidth}{!}{%
        \begin{tikzpicture}[
            % Pfeile und Linien
            arrowstyle/.style={-{Latex[length=2.5mm, width=1.5mm]}, very thick, draw=bshaupt},
            bsverbindung/.style={draw=bshaupt, very thick},
            nfverbindung/.style={draw=bshaupt!60, thick},
            % Gesamtfunktion
            gesamtbox/.style={
                draw=bshaupt, fill=bshaupt, text=white, very thick,
                rounded corners=3pt,
                minimum width=5cm, minimum height=1.8cm,
                align=center, font=\bfseries\large
            },
            % Hauptfunktionen
            hfbox/.style={
                draw=bshaupt, fill=bshell, thick,
                rounded corners=3pt,
                text width=2.4cm, minimum height=2.2cm,
                align=center, font=\footnotesize
            },
            % Nebenfunktionen
            nfblock/.style={
                draw=bshaupt!60, thick, rounded corners=3pt,
                inner sep=4pt,
                text width=2.4cm,
                align=left, anchor=north,
                font=\scriptsize
            }
        ]

            % ==========================================
            % UNTERER TEIL: HAUPTFUNKTIONEN
            % ==========================================

            \node[hfbox] (hf1) at (0,0) {
                \textcolor{bshaupt}{\textbf{Hauptfunktion 1}}\\[1mm]
                Energie-\\versorgung\\(24\,V)
            };

            \node[hfbox, right=0.35cm of hf1] (hf2) {
                \textcolor{bshaupt}{\textbf{Hauptfunktion 2}}\\[1mm]
                Einlesen von seriellen\\Befehlen
            };

            \node[hfbox, right=0.35cm of hf2] (hf3) {
                \textcolor{bshaupt}{\textbf{Hauptfunktion 3}}\\[1mm]
                kinematische\\Berechnungen\\ausführen
            };

            \node[hfbox, right=0.35cm of hf3] (hf4) {
                \textcolor{bshaupt}{\textbf{Hauptfunktion 4}}\\[1mm]
                Ansteuerung\\Bewegung der\\Schritt- und\\Servomotoren
            };

            \node[hfbox, right=0.35cm of hf4] (hf5) {
                \textcolor{bshaupt}{\textbf{Hauptfunktion 5}}\\[1mm]
                Referenzieren
            };

            \node[hfbox, right=0.35cm of hf5] (hf6) {
                \textcolor{bshaupt}{\textbf{Hauptfunktion 6}}\\[1mm]
                Platine und\\Komponenten tragen\\und schützen
            };

            % ==========================================
            % NEBENFUNKTIONEN
            % ==========================================

            \node[nfblock, below=0.6cm of hf1] (nf1) {
                \textbf{Nebenfunktion 1.1:}\\
                24\,V zu 5\,V\\
                Spannungswandlung\par\vspace{3.5mm}
                \textbf{Nebenfunktion 1.2:}\\
                24\,V und 5\,V\\
                Spannungsversorgung\\
                stützen\par\vspace{3.5mm}
                \textbf{Nebenfunktion 1.3:}\\
                24\,V\\
                Treiberversorgung\\
                5\,V ESP/Servo-\\
                Versorgung
            };

            \node[nfblock, below=0.6cm of hf2] (nf2) {
                \textbf{Nebenfunktion 2.1:}\\
                Bluetooth-Serial\\
                Befehle empfangen\par\vspace{3.5mm}
                \textbf{Nebenfunktion 2.2:}\\
                Befehle parsen
            };

            \node[nfblock, below=0.6cm of hf3] (nf3) {
                \textbf{Nebenfunktion 3.1:}\\
                Erreichbarkeit der\\
                Koordinaten prüfen\par\vspace{3.5mm}
                \textbf{Nebenfunktion 3.2:}\\
                inverse Kinematik\\
                lösen, Winkel in\\
                Motorschritte\\
                umrechnen
            };

            \node[nfblock, below=0.6cm of hf4] (nf4) {
                \textbf{Nebenfunktion 4.1:}\\
                Ausgabe von STEP/\\
                DIR-Signalen für die\\
                Schrittmotoren\par\vspace{3.5mm}
                \textbf{Nebenfunktion 4.2:}\\
                Ausgabe PWM-Signal\\
                Servo Motor\par\vspace{3.5mm}
                \textbf{Nebenfunktion 4.3:}\\
                Konfiguration von\\
                Geschwindigkeit und\\
                Beschleunigung
            };

            \node[nfblock, below=0.6cm of hf5] (nf5) {
                \textbf{Nebenfunktion 5.1:}\\
                Auswerten der\\
                Motorströme bei\\
                Anfahren einer\\
                Endlage\par\vspace{3.5mm}
                \textbf{Nebenfunktion 5.2:}\\
                Offset setzen und in\\
                Nullposition fahren
            };

            \node[nfblock, below=0.6cm of hf6] (nf6) {
                \textbf{Nebenfunktion 6.1:}\\
                Platine befestigen\par\vspace{3.5mm}
                \textbf{Nebenfunktion 6.2:}\\
                Komponenten auf der\\
                Platine befestigen\par\vspace{3.5mm}
                \textbf{Nebenfunktion 6.3:}\\
                Wärme abführen\par\vspace{3.5mm}
                \textbf{Nebenfunktion 6.4:}\\
                Zugänglichkeit der\\
                Kabelanschlüsse\\
                gewährleisten
            };

            % Verbindungen Hauptfunktion -> Nebenfunktionen
            \foreach \i in {1,...,6}{
                \draw[nfverbindung] (hf\i.south) -- (nf\i.north);
            }

            % ==========================================
            % OBERER TEIL: GESAMTFUNKTION (Black-Box)
            % ==========================================

            \coordinate (centerHF) at ($(hf1.west)!0.5!(hf6.east)$);
            \node[gesamtbox] (gesamt) at ([yshift=4cm]centerHF) {Gesamtfunktion};

            % Verbindung Gesamtfunktion -> Hauptfunktionen
            \draw[bsverbindung] (gesamt.south) -- ++(0,-0.9) coordinate (bus);
            \foreach \i in {1,...,6}{
                \draw[bsverbindung] (bus) -| (hf\i.north);
            }

            % Eingangspfeil
            \draw[arrowstyle] ([xshift=-6cm]gesamt.west) -- (gesamt.west)
                node[midway, above=2mm, align=center, font=\footnotesize] {
                    \textbf{Eingang}\\
                    Spannung, Steuersignale, Winkel,\\
                    Koordinaten, Geschwindigkeit,\\
                    Beschleunigung
                };

            % Ausgangspfeil
            \draw[arrowstyle] (gesamt.east) -- ([xshift=6cm]gesamt.east)
                node[midway, above=2mm, align=center, font=\footnotesize] {
                    \textbf{Ausgang}\\
                    Rotation der Motoren
                };

        \end{tikzpicture}
    }
    \caption{Funktionsstruktur des Antriebssystems}
    \label{fig:funktionsstruktur}
\end{figure}

\end{document}